% DECISIONES (~1 min cada una). Hipótesis -> evidencia -> resultado -> decisión.
% Fuentes: EXP-002, EXP-003, EXP-005, EXP-006, EXP-007, EXP-008..EXP-010 (LOGBOOK.md).

\begin{frame}{Normalizar la entrada y parar con la validación}
  \begin{columns}[c,onlytextwidth]
    \begin{column}{0.4\textwidth}
      \begin{itemize}
        \item \textbf{Hipótesis}: píxeles en $[0, 255]$ saturan la softmax
        \item \textbf{Evidencia}: \texttt{Rescaling(1/255)}: +\res{gain-normalisation}~pp
        \item \textbf{Resultado}: callbacks sobre el mean recall: la lr baja al estancarse y se para a tiempo
      \end{itemize}
      \vspace{0.2cm}
      \decisionblock{EarlyStopping (paciencia \res{final-es-patience}) y se conserva la mejor época.}
    \end{column}
    \begin{column}{0.55\textwidth}
      \begin{figure}
        \centering
        \safegraphics[width=\textwidth]{figures/stopping.png}{Punto de parada}
        \caption{\footnotesize \figcaption{stopping}}
      \end{figure}
    \end{column}
  \end{columns}
  \speakernote{1}{Vigilar el mean recall dio algo más que vigilar la pérdida o la accuracy, y es la métrica de decisión. La figura es del modelo final: cada semilla para \res{final-es-patience} épocas después de su mejor época.}
\end{frame}

\begin{frame}{Reducir la resolución de entrada}
  \begin{columns}[c,onlytextwidth]
    \begin{column}{0.4\textwidth}
      \begin{itemize}
        \item \textbf{Hipótesis}: con tantas entradas, sobra varianza
        \item \textbf{Evidencia}: a \res{res224-res} píxeles también cae train: casi todas las ReLU mueren
        \item \textbf{Resultado}: \res{phase2-res}~$\times$~\res{phase2-res}: +\res{gain-resolution}~pp, aun con la mejor lr de cada tamaño
      \end{itemize}
      \vspace{0.2cm}
      \decisionblock{Entrada de \res{phase2-res}~$\times$~\res{phase2-res}: el mayor salto.}
    \end{column}
    \begin{column}{0.55\textwidth}
      \begin{figure}
        \centering
        \safegraphics[width=\textwidth]{figures/resolution_sweep.png}{Barrido de resolución}
        \caption{\footnotesize \figcaption{resolution_sweep}}
      \end{figure}
    \end{column}
  \end{columns}
  \speakernote{1}{La caída con la resolución es de optimización, no de sobreajuste.}
\end{frame}

\begin{frame}{Arquitectura: profundidad antes que anchura}
  \begin{columns}[c,onlytextwidth]
    \begin{column}{0.4\textwidth}
      \begin{itemize}
        \item \textbf{Hipótesis}: más capacidad, mejor validación
        \item \textbf{Evidencia}: train sube hasta \res{big3layer-train-acc}; validación se estanca
        \item \textbf{Resultado}: \res{phase3-hidden} empata con la red mayor (\res{phase3-mean-recall}); LeakyReLU evita ReLU muertas
      \end{itemize}
      \vspace{0.2cm}
      \decisionblock{\res{phase3-hidden} con LeakyReLU (\res{phase3-params} parámetros).}
    \end{column}
    \begin{column}{0.55\textwidth}
      \begin{figure}
        \centering
        \safegraphics[width=\textwidth]{figures/capacity.png}{Curva de capacidad}
        \caption{\footnotesize \figcaption{capacity}}
      \end{figure}
    \end{column}
  \end{columns}
  \speakernote{1}{Leer la curva de capacidad: la separación entre train y validación crece con los parámetros.}
\end{frame}

\begin{frame}{Optimizador y parámetros de entrenamiento}
  \begin{columns}[c,onlytextwidth]
    \begin{column}{0.4\textwidth}
      \begin{itemize}
        \item \textbf{Hipótesis}: la learning rate es el hiperparámetro clave
        \item \textbf{Evidencia}: Adam casi insensible a la lr; SGD + momentum con un óptimo marcado
        \item \textbf{Resultado}: empatan dentro del ruido (\res{seed-noise-pp}~pp entre semillas)
      \end{itemize}
      \vspace{0.2cm}
      \decisionblock{Adam, $\beta_1$ = \res{final-beta1}, batch \res{final-batch}: épocas \res{speedup-batch256} veces más rápidas.}
    \end{column}
    \begin{column}{0.55\textwidth}
      \begin{figure}
        \centering
        \safegraphics[width=\textwidth]{figures/lr_sweep.png}{Barrido de learning rate}
        \caption{\footnotesize \figcaption{lr_sweep}}
      \end{figure}
    \end{column}
  \end{columns}
  \speakernote{1}{El batch grande cuesta unos \res{cost-batch256}~pp frente a batch \res{phase4-batch}, pero permite decidir con \res{final-seeds} semillas. Ningún schedule mejoró a ReduceLROnPlateau.}
\end{frame}
